\subsection*{4.4 Proofs of Theorem3 and Corollary4}
With the machinery developed in the previous section, and using in a black-box manner, the construction of axis-evasive partitions we obtain in Section5, we are finally ready to prove Theorem3 and Corollary4. We start by giving a more formal statement of Corollary4 (Theorem46, stated at the start of this item).
\paragraph{A remark regarding the error}
Note that there is another implicit constraint on $\rho_{\mathrm{in}}$ and $\alpha$ that originates from the error. Indeed, to make the result non-trivial, one must have $\varepsilon_{\mathrm{out}}<1$ which, in turn, implies some bound on $\ell$ and then, through Equation(4.12), a constraint on $\rho_{\mathrm{in}}$ and $\alpha$. However, if that turns out to be a problem for the regime of parameters one is interested in, the probability to output $\bot$ can be reduced by repetition. Thus, by performing an alternating sequence of such error (or failure) reductions and rate amplifications, one can resolve this issue. Note that unlike for LDC, the error reduction has no cost in query complexity, and it certainly has no effect on the smoothness nor on the rate. It does, however, effects the running-time.
As mentioned above, our proof relies on an explicit axis-evasive partition that we construct in Section5. Formally,
\setcounter{nttheorem}{46}
\begin{nttheorem}[Author Theorem47]
Let $P$ be a set of size $q$, where $q$ is an odd prime power. Let $c$ be an even integer such that $c+1\mid q+1$, and $c\leq\sqrt q/10$. Then, there exists a $(c,4c^2)$-axis evasive partition of $P^2$ with at most $2q^2/(c+1)$ parts.
\end{nttheorem}
Our proof of Theorem46 is done by applying Proposition45 several times, iteratively, where in each iteration we square the size of the set $P$ obtained by the previous iterative step. Note, however, that Theorem47 requires the set size $|P|$ to be an odd prime power $q$ with the property that $c+1\mid q+1$. It is best to choose $c$ the same in all applications of Proposition45. However, note that if we start an iteration with a set of size $q$ and so end the iteration with a set of size $q^2$ then the condition will fail to hold at the beginning of the following iteration. Indeed if $c+1\mid q+1$ then $q\equiv-1\pmod{c+1}$ and so $q^2\equiv1\pmod{c+1}$. To overcome this technicality, we do not work with the set obtained by the previous iteration as is. Instead, we find a prime--not much smaller than $q^2$--that has the desired residue $-1$ modulo $c+1$. To this end we rely on the Siegel--Walfisz Theorem~\cite{ref34,ref37} which refines Dirichlet's theorem on primes in arithmetic progressions. The state the Siegel--Walfisz Theorem we set some notation. For an integer $m\geq1$, we denote Euler's totient function, that counts the positive integers up to $m$ that are relatively prime to $m$, by $\varphi(m)$. For integers $n,m,r$, we denote the number of (positive) primes less than or equal to $n$ which are congruent to $r$ modulo $m$ by $\pi(n;m,r)$. The Eulerian logarithmic integral is given by $\operatorname{Li}(x)=\int_2^x dt/\ln t$.
\begin{nttheorem}[Author Theorem48, {\cite{ref34,ref37}}]
For every constant $e\geq1$ there exists a constant $c=c(e)$ such that the following holds. Let $n,m,r$ be positive integers such that $m\leq(\log n)^e$, and $m,r$ coprimes. Then,
\[\left|\pi(n;m,r)-\frac{\operatorname{Li}(n)}{\varphi(m)}\right|=O\left(n\cdot2^{-c\sqrt{\log n}}\right).\]
\end{nttheorem}
We have the following straightforward corollary.
\begin{nttheorem}[Author Corollary49]
For every constant $e\geq1$ there exist constants $c=c(e),n_0=n_0(e)$ such that the following holds. Let $m,r$ be coprime integers, $m>0$. Let $n\geq n_0$ be an integer such that $m\leq(\log n)^e$. Then, there exists a prime $p\in[n-\Delta,n]$, where $\Delta=cn/\log n$, such that $p\equiv r\pmod m$.
\end{nttheorem}
\begin{proof}
To prove the corollary, it suffices to show that $\pi(n;m,r)>\pi(n-\Delta;m,r)$. By Theorem48, there exist constants $n_0,c'$ such that for every $n\geq n_0$,
\[\left|\pi(n;m,r)-\frac{\operatorname{Li}(n)}{\varphi(m)}\right|\leq c'n\cdot2^{-c\sqrt{\log n}}.\]
Thus, it suffices to show that
\[\frac{\operatorname{Li}(n)}{\varphi(m)}-c'n\cdot2^{-c\sqrt{\log n}}>\frac{\operatorname{Li}(n-\Delta)}{\varphi(m)}+c'(n-\Delta)\cdot2^{-c\sqrt{\log(n-\Delta)}}.\]
As we may assume that $\Delta\leq n/2$, it suffices to prove that
\begin{equation}\operatorname{Li}(n)-\operatorname{Li}(n-\Delta)\geq2c'\varphi(m)n\cdot2^{-c\sqrt{\log(n/2)}}.\tag{4.13}\end{equation}
It is well-known that
\[\operatorname{Li}(x)=c_1+\frac x{\ln x}+O\left(\frac x{\ln^2x}\right),\]
where $c_1=\int_{t=0}^2dt/\ln t$ is some constant. Therefore,
\[\operatorname{Li}(n)-\operatorname{Li}(n-\Delta)\geq\frac\Delta{\ln(n/2)}-\frac{c''n}{\ln^2n}.\]
for some constant $c''$. By our assumption on $\Delta$ we can choose the parameter $c$ in the definition of $\Delta$ such that the right hand side is bounded below by $n/\ln^2n$. The proof then follows by Equation(4.13) and noting that $\varphi(m)\leq m\leq(\log n)^e=o(2^{-c\sqrt{\log(n/2)}})$.
\end{proof}
We turn to formally define and analyze the operation of projecting a dual SLR over $\mathbb F^P$ on a (large) subset of $P$.
\begin{ntdefinition}[Author Definition50]
Let $P$ a set and $P'\subseteq P$. Let $p'\in P'$ and $D$ be a distribution with $\operatorname{supp}(D)\subseteq F_{p'}\cup\{\bot\}$. We define the $D|_{P'}$ as follows: To sample from $D|_{P'}$, sample $f\sim D$. If $f=\bot$, output $\bot$; if $f\in F_{p'}$, output $f|_{P'}$. We refer to $D|_{P'}$ as the distribution $D$ projected to $P'$.
\end{ntdefinition}
\begin{ntdefinition}[Author Definition51]
Let $P$ be a set, $\mathbb F$ a field. Let $D=\{D_p\mid p\in P\}$ be a collection of distributions, where for each $p\in P$, $\operatorname{supp}(D_p)\subseteq F_p\cup\{\bot\}$. Let $P'\subseteq P$. We define $D|_{P'}$ to be the collection $D$ projected to $P'$, that is, $D|_{P'}=\{D_{p'}|_{P'}\mid p'\in P'\}$.
\end{ntdefinition}
\begin{ntdefinition}[Author Definition52]
Let $P$ be a set, $\mathbb F$ a field and let $L$ be a linear subspace of $\mathbb F^P$. Let $P'\subseteq P$. We denote by $L|_{P'}$ the linear subspace $L$ projected to $P'$, namely, $L|_{P'}=\{f|_{P'}\mid f\in L\}$.
\end{ntdefinition}
\begin{ntclaim}[Author Claim53]
Let $P$ be a set, $\mathbb F$ a field, $(D,L)$ a $(q,\tau,\varepsilon,\rho)$-dual SLR over $\mathbb F^P$, and let $P'\subseteq P$. Then, $(D|_{P'},L_{P'})$ is a $(q,\tau,\varepsilon,\rho')$-dual SLR over $\mathbb F^{P'}$, where $\rho'=1-\frac{|P|}{|P'|}(1-\rho)$.
\end{ntclaim}
\begin{proof}
That the smoothness $\tau$, as well as $q$ and $\varepsilon$, all stay the same after projecting the dual SLR to $P'$, follows immediately from the definitions. The assertion regarding the rate of the induced SLR, $\rho'$, readily follows as we have that
\[\dim(L|_{P'})\leq\dim(L)\leq(1-\rho)|P|=\left(1-\left(1-\frac{|P|}{|P'|}(1-\rho)\right)\right)|P'|.\]
\end{proof}
\begin{ntclaim}[Author Claim54]
There exists a universal constant $m_0$ such that the following holds. Let $P$ be a set of size $m\geq m_0$. Let $\mathbb F$ be a field, and let $(D^P,L^P)$ be a $(q_{\mathrm{in}},\tau_{\mathrm{in}},\varepsilon_{\mathrm{in}},\rho_{\mathrm{in}})$-dual SLR over $\mathbb F^P$. Let $c\leq\log m$ be an integer. Then, there exists a set $P'$ of size
\[|P'|\geq\left(1-O\left(\frac1{\log m}\right)\right)m^2,\]
and a $(q_{\mathrm{out}},\tau_{\mathrm{out}},\varepsilon_{\mathrm{out}},\rho_{\mathrm{out}})$-dual SLR $(D^{P'},L^{P'})$ over $\mathbb F^{P'}$, where
\begin{align*}q_{\mathrm{out}}&\leq2cq_{\mathrm{in}}^3,\\\tau_{\mathrm{out}}&\leq40c^3q_{\mathrm{in}}\tau_{\mathrm{in}}^2,\\\varepsilon_{\mathrm{out}}&\leq80c^3q_{\mathrm{in}}(\tau_{\mathrm{in}}^2+\varepsilon_{\mathrm{in}}),\\\rho_{\mathrm{out}}&\geq1-(1-\rho_{\mathrm{in}})^2-O(1/c).\end{align*}
\end{ntclaim}
\begin{proof}
By Corollary49 applied with $n,m,r$ in the notation of Corollary49 set to $m,c+1,-1$ in the notation of this claim, respectively, there exists some prime $p\leq m$ such that $m-p=O(m/\log m)$, and $c+1\mid p+1$. Take $P'$ to be an arbitrary subset of $P$ of size $p$. By Claim53, $(D|_{P'},L|_{P'})$ is a $(q_{\mathrm{in}},\tau_{\mathrm{in}},\varepsilon_{\mathrm{in}},\rho')$-dual SLR on $P'$, where
\[\rho'=1-\frac mp(1-\rho_{\mathrm{in}})\geq\rho_{\mathrm{in}}-O\left(\frac1{\log m}\right).\]
By Theorem47 applied to $P'$, which observe is indeed applicable as $c+1\mid p+1$, there exists an explicit $(c,4c^2)$-axis evasive partition $\mathcal R$ of $(P')^2$ with at most $t=2p^2/(c+1)$ parts. With that partition, we can now apply Proposition45 to $(D|_{P'},L|_{P'})$ and get that $(D|_{P'},L|_{P'})^2_{\mathcal R}$ is a $(q_{\mathrm{out}},\tau_{\mathrm{out}},\varepsilon_{\mathrm{out}},\rho_{\mathrm{out}})$-dual SLR with the stated parameters. Note that the assertion regarding the rate follows as $c\leq\log m$,
\end{proof}
The following proposition is a more formal and accurate restatement of Theorem3.
\begin{ntproposition}[Author Proposition55]
There exist universal constants $0<c'<1$ and $c'',m',\ell'\geq1$ such that the following holds. Let $P$ be a set of size $m\geq m'$. Let $\mathbb F$ be a field, and let $(D^P,L^P)$ be a $(q_{\mathrm{in}},\tau_{\mathrm{in}},\varepsilon_{\mathrm{in}},\rho_{\mathrm{in}})$-dual SLR over $\mathbb F^P$. Let $\ell=2^r$ for an integer $r\geq1$, and assume that $\ell\geq\ell'$. Let $c$ be an integer such that $c''\ell^2\leq c\leq c'\log m$. Then, there exists a set $P_\ell$ of size $m^\ell/2\leq|P_\ell|\leq m^\ell$, and a $(q_\ell,\tau_\ell,\varepsilon_\ell,\rho_\ell)$-dual SLR $(D^{P_\ell},L^{P_\ell})$ over $\mathbb F^{P_\ell}$, where
\begin{align*}q_\ell&\leq(2cq_{\mathrm{in}})^{\ell^{\log3}},\\\tau_\ell&=O((c^3q_{\mathrm{in}})^{\ell^{\log3}})\cdot\tau_{\mathrm{in}}^\ell,\\\varepsilon_\ell&\leq O((c^4q_{\mathrm{in}})^{\ell^{\log3}})\cdot(\tau_{\mathrm{in}}+\varepsilon_{\mathrm{in}}),\\\rho_\ell&\geq1-(1-\rho_{\mathrm{in}})^\ell-O\left(\frac{\ell^2}c\right),\end{align*}
where, recall, the $\log$ function is taken base $2$.
\end{ntproposition}
\begin{proof}
We construct a sequence of $(q_t,\tau_t,\varepsilon_t,\rho_t)$-dual SLR $(D^{P_t},L^{P_t})$ for $t=0,1,\ldots,r=\log\ell$, and show that the last dual-SLR in the sequence has the stated parameters. The first dual-SLR, $(D^{P_0},L^{P_0})$, is taken to be the $(q_{\mathrm{in}},\tau_{\mathrm{in}},\varepsilon_{\mathrm{in}},\rho_{\mathrm{in}})$-dual SLR $(D^P,L^P)$ that is given by the hypothesis of the proposition. After constructing $(D^{P_t},L^{P_t})$, we obtain $(D^{P_{t+1}},L^{P_{t+1}})$ by applying Claim54 to $(D^{P_t},L^{P_t})$ with the parameter $c$ taken to be $c$ from the statement of this proposition. Note that, as required by the claim, $c\leq\log m$. Note that, by taking $m'$ to be a large enough constant, all other dual SLR in the sequence will have $|P_t|\geq m$ as well, and so we can apply Claim54 to them. Denote $m_t=|P_t|$. We begin by bounding $m_t$ from below. Indeed, by Claim54, and using that $1-x\geq e^{-2x}$ for $x\leq1/2$, we can pick the constant $c''$ such that
\[m_t\geq e^{-c''/\log m_{t-1}}m_{t-1}^2\geq e^{-c''/\log m_0}m_{t-1}^2,\]
where the last inequality follows as, for a large enough constant $m'$, the sequence $(m_t)_t$ is monotone increasing. We invoke Claim82 with $a=e^{c''/(2\log m_0)}$ and $b=2$ to conclude that
\[m_t\geq m_0^{2^t}e^{-c''2^t/\log m_0}\geq\tfrac12m_0^{2^t},\]
where the last inequality follows as $t\leq r=\log\ell$ and, recall, we take $\ell\leq c'\log m$ for a sufficiently small constant $c'>0$. In particular, $m_r\geq m^\ell/2$ as stated.
By Claim54, for every $t\geq1$ we have $q_t\leq2cq_{t-1}^3$. It is straightforward to prove by that
\begin{equation}q_t\leq(2cq_{\mathrm{in}})^{3^t},\tag{4.14}\end{equation}
which readily implies the assertion regarding the query complexity. We turn to analyze the rate. Denote $\beta_t=1-\rho_t$. Claim54 implies that $\beta_t\leq\beta_{t-1}^2+c'''/c$, for some constant $c'''>0$. By induction on $t$, we get that $\beta_t\leq\beta_0^{2^t}+c'''4^t/c$. Indeed, the base case $t=0$ is obvious. Now, by the induction hypothesis,
\[\beta_t\leq\beta_{t-1}^2+\frac{c'''}c\leq\left(\beta_0^{2^{t-1}}+\frac{c'''}c4^{t-1}\right)^2+\frac{c'''}c.\]
One can easily verify that the right hand side is bounded above by the desired bound $\beta_0^{2^t}+c'''4^t/c$ provided that $2^tc'''/c\leq1$. As $t\leq r$ and $2^r=\ell$, the latter inequality follows assuming $c'''\ell\leq c$. As we assume $c\geq c''\ell^2$, it suffices to choose $\ell'$ from the statement of the proposition to be a constant larger than the constant $c'''/c''$. We conclude that,
\[\beta_r\leq\beta_0^\ell+O\left(\frac{4^r}c\right)=\beta_0^\ell+O\left(\frac{\ell^2}c\right),\]
which implies the assertion regarding the rate.
As for the smoothness, by Claim54, and using Equation(4.14), we have that
\[\tau_t\leq40c^3q_{t-1}\tau_{t-1}^2\leq40c^3(2cq_{\mathrm{in}})^{3^{t-1}}\tau_{t-1}^2,\]
from which it is easy to verify that
\[\tau_t\leq(40c^3)^{2^t}(2cq_{\mathrm{in}})^{3^t}\tau_0^{2^t},\]
and the assertion regarding the smoothness readily follows. Last is the error which we leave to the reader to verify.
\end{proof}
We can now easily deduce Theorem46.
\begin{proof}[Author proof of Theorem46]
The proof readily follows from Proposition55 by taking $\ell$ as defined in Equation(4.12), and with $c$ in the notation of Proposition55 taken to be $c=\Theta(\ell^2/\alpha)$. Note that this choice of parameters satisfies the hypothesis of Proposition55 as indeed implied by Equation(4.11) and by taking $c'$ to be a sufficiently large constant. It is easy to verify that the rate is $1-\alpha$ with our choice of $c,\ell$, and the remaining assertions readily follow by Proposition55.
\end{proof}
\setcounter{nttheorem}{81}
\begin{ntclaim}[Author Claim82, used above; p1:48]
Let $(n_t)_{t\in\mathbb N}$ be a sequence of positive integers such that $n_t\geq(n_{t-1}/a)^b$ for some $a,b>1$. Then, for every $t\geq1$ we have that $n_t\geq(n_0/a^{h(b,t)})^{b^t}$, where $h(b,t)=\sum_{i=0}^{t-1}1/b^i$.
\end{ntclaim}
One can prove the following easy claim by induction (author's supplied justification for Claim82).
